% CHECKED extension: complete self-reading and two independent refinement audits passed.
\subsection{A weaker remaining arithmetic condition}
\label{sec:pole-nonvanishing}

The proof of Theorem~\ref{thm:pole-general} uses the unit cup coefficient only to ensure that a selected numerator stays nonzero on a pole divisor. This gives a weaker sufficient condition, expressed in the actual quadratic companion errors.

Retain the construction of $H$, its compact basis, $q_e(z)$, and the integral numerator functions $C_{e;a,b}(T_1,T_2)$ from the preceding proof. That construction, up to the assertion (G.3), uses Vandiver but no cup-surjectivity hypothesis. Let $L/\mathbb Q_p$ split the finitely many polynomials $q_e$, let $\mathcal O_L$ be its valuation ring, and let $\pi$ be a uniformizer. Every root $\alpha$ lies in $1+\pi\mathcal O_L$.

\begin{theorem}[K17-POLE-NV]\label{thm:pole-nv}
Let $p$ be an odd prime satisfying Vandiver. Suppose that for every occurring even character $e$ and every distinct root $\alpha$ of $q_e$, there is an ordered pair of odd characters $(a,b)$, with $a+b\equiv e\pmod{p-1}$, such that
\[
 C_{e;a,b}\!\left(T,\frac{\alpha}{1+T}-1\right)\ne0
 \quad\text{in }\mathcal O_L[[T]][1/\pi].                    \tag{NV}
\]
Then $\Omega_p^*=\Omega_p$. The pair $(a,b)$ may depend on $\alpha$, and no assumption on the multiplicities of the roots is needed.
\end{theorem}

Equivalently, for each $(e,\alpha)$ at least one matching numerator is not divisible by
\[
 D_\alpha=(1+T_1)(1+T_2)-\alpha
\]
in $\mathcal O_L[[T_1,T_2]][1/\pi]$. The equivalence follows from the formal change of variables $S=D_\alpha$ and $T_2=(\alpha+S)/(1+T_1)-1$; its constant displacement is topologically nilpotent. Thus the condition concerns restriction to the divisor, rather than just nonvanishing somewhere in the two-variable domain.

\begin{proof}
All the arithmetic module, compact free basis, actual rational realization, companion identity, and partial-fraction steps of the preceding proof remain valid without (G.3). They give degree-two formal letters with leading target kernels
\[
 \frac{C_{e;a,b}(T_1,T_2)}{D_\alpha^{\,k}},
 \qquad 1\le k\le n_{e,\alpha}.
\]
The finite scalar extension and homogeneous change of letters preserve the bounded-denominator coefficient measures and the first nonzero assigned Magnus degree. Restrict the test weights to the same sufficiently large Cartesian tails to avoid all exceptional denominator zeros.

Suppose a nontrivial kernel element has first assigned degree $d$. Among its nonzero colored tilings choose one with the maximum number of pairs. Fix its pair positions and $(e,\alpha)$ colors, and then choose a componentwise maximal pole-order vector $(k_j)$. At each chosen pair use the odd target types supplied by (NV) for that particular $(e_j,\alpha_j)$; at each singleton retain its specified source orbit type.

The vanishing target coefficients and the successive analytic identity theorem give exactly (G.9) in the bounded-denominator fraction ring. At the disjoint selected pairs put $S_j=D_{j,\alpha_j}$ and retain the unpaired variables and one left variable from each pair as base variables. Let $K_0$ be the fraction field of the bounded-denominator power-series ring in these base variables. The same denominator classification embeds (G.9) in
\[
 K_0[[S_j:j\in J]][S_j^{-1}:j\in J].
\]
Other-root denominators at a selected pair have nonzero constant in $K_0$; denominators at different positions connect different independent blocks and also have nonzero constant. Wrong-character numerators vanish identically, as already proved using the Cartesian weight grids.

Extract the coefficient of $\prod_{j\in J}S_j^{-k_j}$. A pattern missing a selected pair or using a different root there cannot supply the corresponding negative exponent. A pattern with additional pairs has zero coefficient by maximality of the pair count. Among patterns with the same positions and colors, only a componentwise greater or equal pole-order vector could contribute, and the chosen maximality excludes every strict such vector. Therefore the coefficient is precisely
\[
 A_{\mathcal P}\prod_{j\in J}
 C_{e_j;a_j,b_j}\!\left(T_j,\frac{\alpha_j}{1+T_j}-1\right).
\]
The selected coefficient $A_{\mathcal P}$ is nonzero. Each other factor is nonzero by (NV), the variables belonging to distinct pairs are independent, and the base power-series ring is a domain. Their product is consequently nonzero in $K_0$, contradicting the extracted zero coefficient. When there are no pairs, this is the already checked odd-moment separation argument.

Thus the actual rational action is injective on $H$. The cyclotomic character detects any kernel outside $H$, and the rational action is induced by the pro-$p$ outer action. This proves the asserted equality.
\end{proof}

\paragraph{Comparison with the cup condition.}
A numerator with unit constant coefficient is a unit in $\mathbb Z_p[[T_1,T_2]]$ and remains a unit under every indicated divisor substitution. Hence the cup condition in Theorem~\ref{thm:pole-general} implies (NV). The converse analytic implication fails: the formal numerator $C=p$ has no nonzero reduction modulo $p$, but restricts to a nonzero constant on every divisor. This is an example of analytic data, not an arithmetic prime. No additional prime satisfying (NV) but failing the cup hypothesis has been established here.

Repeated roots explain why the divisor qualification matters. A globally nonzero numerator divisible by $D_\alpha$ cancels a highest pole order; it need not suffice for the present separation argument. Condition (NV) rules out such cancellation on at least one odd pair for each $(e,\alpha)$.

\paragraph{Ordinary basis choices.}
For a fixed $\gamma$, replacing a cyclic lift $y_e$ by another lift of the same abelian class changes its length-two target coefficients by a bounded analytic commutator transform $A$. Consequently
\[
 B'_e=B_e+A,\qquad C'_e=C_e+q_e(z)A,
 \qquad z=(1+T_1)(1+T_2),
\]
so their restrictions at $z=\alpha$ agree. Replacing the cyclic generator by $b_e(\gamma)v_e$, where $b_e$ is a unit modulo $q_e$, instead gives
\[
 B'_e=b_e(z)B_e+A,\qquad
 C'_e=b_e(z)C_e+q_e(z)A.
\]
The restricted numerators differ by $b_e(\alpha)\ne0$. Changing an odd $\Lambda$ basis multiplies its normalized linear moments by a power-series unit; this divides the corresponding normalized numerator by unit factors in its two variables and again preserves (NV). Scalar extension preserves nonzero power series. These facts make the condition independent of these routine choices; the cyclotomic generator may simply remain fixed by $\chi(\gamma)=1+p$.

\paragraph{Arithmetic scope of the refinement.}
The criterion records exactly where the separation proof uses its
quadratic input. The next subsection supplies the stronger cup
condition at every Vandiver prime, so (NV) introduces no independent
arithmetic hypothesis in that range. A failure of (NV) would not by
itself prove a nontrivial kernel: a source direction might still be
detected at a larger target generator length. An argument beyond
Vandiver would also have to accommodate the different module structure.
